\begin{minipage}{0.72\linewidth}
    Deux électrodes, de surface $S_1=\qty{2,0}{\square\cm}$ et distantes de
    $L_1=\qty{2,0}{\cm}$ l'une de l'autre plongent dans une solution ionique
    contenant un seul soluté.
    
    On réalise le montage représenté sur le schéma ci-dessous~; le voltmètre
    indique une tension $U_1=\qty{1,0}{\volt}$ et l'ampèremètre une intensité
    $I_1=\qty{5,0}{\mA}$.

    Après avoir rapproché les électrodes, on note une tension de
    $\qty{1,0}{\volt}$ et une intensité de $\qty{7,0}{\mA}$.
    \begin{enumerate}[series=x]
        \item Calculer la conductance $G$ de la solution dans les deux cas.
        \item Calculer la distance entre les deux électrodes dans le deuxième
              cas.
        \item Calculer la constante de cellule dans les deux cas (préciser
              l'unité).
        \item Calculer la conductivité de la solution.
    \end{enumerate}
\end{minipage}
\hfill
\begin{minipage}{0.25\linewidth}
    \flushright
    \begin{tikzpicture}
        \newdimen\W \W=3cm
        \newdimen\H \H=2.5cm
        \newdimen\d \d=1.5cm
        \pgfmathparse{0.6\W}\let\L=\pgfmathresult
        \node[inner sep=0pt,fill=lightgray,minimum width=\W,
            minimum height=0.7\H] (S) {};
        \node[inner sep=0pt,minimum width=\W,
            minimum height=\H,anchor=south] (B) at (S.south) {};
        \draw[draw,thick] (B.north east) -- (B.south east) -| (B.north west);
        \node[inner sep=0pt,minimum width=\L,
            minimum height=0.55\H,anchor=south] (C) at (B.south) {};
        \foreach \a in { east, west}
        \node[anchor=north,trapezium,draw,fill=copper,trapezium left angle=120,
            trapezium right angle=60,shape border rotate=90,minimum width=7mm,
            minimum height=6mm,inner sep=0pt] at (C.north \a) {$S$};
        \draw (C.north east) to[short] ++(0,\d)
            to[short] ++(\d,0)
            to[short] ++(0,1.5\d)
            to[rmeter,t={\footnotesize GBF}] ++(-\d-0.5*\L,0)
            to[rmeter,t=A] ++(-\d-0.5*\L,0) coordinate(t1)
            to[short] ++(0,-1.5\d) coordinate(t2)
            to[short] ++(\d,0)
            to[short] (C.north west);
        \draw ($(t1)!0.5!(t2)$) to[rmeter,t=V] ++(2\d+\L,0);
        \draw[<->,shorten <=1pt,shorten >=1pt]
            ([yshift=1cm]C.north west) -- node[above] {$L$} ++(\L pt,0);
    \end{tikzpicture}
\end{minipage}
\begin{enumerate}[resume=x]
    \item Dans une troisième expérience, les électrodes sont identiques à celles
          de la première expérience (même distance $L_1$ et même surface $S_1$),
          la solution ionique contient le même soluté et on opère à la même
          température; le voltmètre indique une tension de $\qty{1,0}{\volt}$ et
          l'ampèremètre une intensité de \qty{15}{\mA}. Quelles sont les
          grandeurs physiques qui ont changé dans cette expérience~?

          Décrire brièvement l'expérience à réaliser pour vérifier les résultats
          précédents.
\end{enumerate}
